\chapter{Frequent Urination}

\section*{Definition}
Frequent urination is needing to pass urine more often than usual. Frequency may occur with normal, increased, or reduced urine volume and should be distinguished from polyuria.

\section*{What Happens in Your Body}
The symptom reflects irritation, inflammation, obstruction, altered secretion, infection, or dysfunction in the relevant organ system. The pattern, duration, associated symptoms, exposures, and examination findings determine whether it is self-limited or requires investigation.

\section*{Causes}
\begin{itemize}
  \item Urinary infection, diabetes, diuretics, pregnancy, overactive bladder, enlarged prostate, excess fluids, caffeine, kidney disease, and electrolyte disorders.
  \item Medication effects, environmental exposures, and chronic disease may contribute.
  \item Less common but important causes include malignancy, vascular disease, or organ failure when symptoms persist or progress.
\end{itemize}

\section*{When to See a Doctor}
Seek urgent evaluation for Inability to urinate, fever with flank pain, visible blood, vomiting, severe weakness, confusion, or new neurologic symptoms. Arrange routine or prompt assessment for persistent, recurrent, unexplained, or function-limiting symptoms.

\section*{Related Symptoms}
\begin{itemize}
  \item Pain, fever, fatigue, nausea, or changes in normal organ function
  \item Bleeding, weight change, shortness of breath, weakness, or neurologic symptoms when a serious cause is present
\end{itemize}

\section*{Self-Care / Home Management}
Avoid known triggers, maintain reasonable hydration, record timing and associated symptoms, and use only treatments appropriate for the suspected cause. Do not use leftover antibiotics or delay urgent care because symptoms temporarily improve.

\section*{How Your Doctor Will Evaluate You}
\begin{itemize}
  \item History addresses onset, duration, severity, triggers, exposures, medications, medical conditions, pregnancy status when relevant, and warning symptoms.
  \item Examination focuses on vital signs and the organ systems suggested by the symptom.
  \item Testing may include blood or urine studies, cultures, imaging, physiologic testing, or specialist examination according to the clinical pattern.
\end{itemize}

\section*{Treatment Options}
Treatment is directed at the cause and may include supportive care, targeted medication, treatment of infection or inflammation, correction of metabolic disease, or specialist procedures. Persistent symptoms require reassessment rather than indefinite symptomatic treatment.

\section*{References}
\begin{thebibliography}{99}
\bibitem{symptom_frequent_urination:ref1} Jameson JL, et al. \textit{Harrison's Principles of Internal Medicine}. 21st ed. McGraw-Hill; 2022.
\bibitem{symptom_frequent_urination:ref2} National Institute for Health and Care Excellence. Relevant assessment and management guidance. NICE; accessed 2025.
\bibitem{symptom_frequent_urination:ref3} Cleveland Clinic. Health Library symptom guidance. Accessed 2026.
\end{thebibliography}
